% Notation and conventions used throughout the appendix; the tolerances of
% the checkers. The verification protocol, the certificate classes and the
% threats to validity (verification.tex) follow as the last subsections of
% this section.
\section{Notation and verification protocol}\label{app:notation}

\subsection{Notation and conventions}\label{app:conventions}

\paragraph{Bilevel problems in BOBILib form.}
The leader controls $x\in\mathbf{R}^{n_x}$ and the follower controls
$y\in\mathbf{R}^{n_y}$. As in BOBILib,\cite{thurauf2026bobilib} the leader minimises:
\begin{align}
  \min_{x,y}\quad & c^{\top}x+d^{\top}y \label{eq:bl}\\
  \text{s.t.}\quad & Ax+By\ge a,\qquad x_i\in\mathbf{Z}\ (i\in I_u),\qquad y\in S(x),\notag\\
  S(x)={}&\operatorname*{arg\,min}_{y'}\bigl\{f^{\top}y' : Cx+Dy'\ge b,\ y'_i\in\mathbf{Z}\ (i\in I_l)\bigr\}. \notag
\end{align}
All data are rational. Variable bounds, equality rows and ranged rows of
the instance files are written as pairs of rows of the form ``$\ge$''; when
an argument depends on the side of a row we say so. For a leader choice $x$
the follower's problem is
$P(x)=\min\{f^{\top}y : Cx+Dy\ge b,\ y_i\in\mathbf{Z}\ (i\in I_l)\}$ and its value
is $\varphi(x)\in\mathbf{R}\cup\{-\infty,+\infty\}$ ($+\infty$ when $P(x)$ is
infeasible). So $S(x)\neq\emptyset$ exactly when $\varphi(x)$ is finite and
attained.

An upper-level row is a \emph{coupling row} if it has a nonzero coefficient
on a follower variable. A leader variable is a \emph{linking variable} if its
column in $C$ is nonzero; $x_L$ denotes the linking variables and $C_L$ their
columns, so $Cx=C_Lx_L$, and $P(x)$, $\varphi(x)$ and $S(x)$ depend on $x$ only
through $x_L$.

\begin{definition}[Bilevel feasible point, optimistic convention]\label{def:bilevel-feasible}
A point $(x,y)$ is \emph{bilevel feasible} if $Ax+By\ge a$, $x_i\in\mathbf{Z}$ for
$i\in I_u$, and $y\in S(x)$. The bilevel optimum $\mathrm{OPT}$ is the infimum
of $c^{\top}x+d^{\top}y$ over bilevel feasible points ($+\infty$ if there is
none).
\end{definition}

Problem~\eqref{eq:bl} minimises over the pair $(x,y)$, so when $S(x)$ has
several elements the follower is assumed to pick the one that is best for the
leader. This is the optimistic convention, which BOBILib uses for all of its
reported values. For a fixed $x$ the optimistic choice solves the
\emph{tie-break problem}
\[
  \min\bigl\{c^{\top}x+d^{\top}y : Ax+By\ge a,\ y\ \text{feasible for}\ P(x),\ f^{\top}y\le\varphi(x)\bigr\}.
\]

\paragraph{High point relaxation.}
The \emph{high point relaxation} (HPR) of~\eqref{eq:bl} drops the condition
$y\in S(x)$ and keeps every row, bound and integrality requirement of both
levels and the leader's objective:
\[
  \mathrm{HPR}=\min\bigl\{c^{\top}x+d^{\top}y : Ax+By\ge a,\ Cx+Dy\ge b,\
  x_i\in\mathbf{Z}\ (i\in I_u),\ y_i\in\mathbf{Z}\ (i\in I_l)\bigr\}.
\]

\paragraph{Packing interdiction.}
Items are $N=\{1,\dots,n\}$. The leader chooses $x\in\{0,1\}^n$ subject to
budget rows $Gx\le g$ with $G\in\mathbf{Z}_{\ge 0}^{m_u\times n}$ and
$g\in\mathbf{Z}_{\ge 0}^{m_u}$. The follower then solves
\[
  \Phi(x)=\max\bigl\{p^{\top}y : Wy\le c,\ y\le \mathbf{1}-x,\ y\in\{0,1\}^n\bigr\}
\]
with profits $p\in\mathbf{Z}_{\ge 0}^n$, weights $W\in\mathbf{Z}_{\ge 0}^{m\times n}$
and capacities $c\in\mathbf{Z}_{\ge 0}^m$, and the leader solves
$\mathrm{OPT}=\min\{\Phi(x) : Gx\le g,\ x\in\{0,1\}^n\}$. In BOBILib form the
leader's objective is $p^{\top}y$ and the follower's is $-p^{\top}y$, so the
game is zero-sum: every $y\in S(x)$ gives the leader the same value $\Phi(x)$,
and the optimistic and pessimistic readings coincide. We identify $x$ with the
set $X=\{j : x_j=1\}$ of interdicted items and write $q(Y)=\sum_{j\in Y}q_j$ for
a vector $q$ and a set $Y\subseteq N$. With
$K(T)=\max\{p(Y): Y\subseteq T,\ \sum_{j\in Y}W_{\cdot j}\le c\}$ the best packing
of the items in $T$, we have $\Phi(X)=K(N\setminus X)$.

KIP (the knapsack interdiction problem) is the case $m_u=m=1$. There we write
$\gamma_j=G_{1j}$ for the cost of interdicting item $j$, $\beta=g_1$ for the
leader's budget, $w_j=W_{1j}$ for the follower's weights and $c=c_1$ for its
capacity, so
$\Phi(X)=\max\{p(Y): Y\subseteq N\setminus X,\ w(Y)\le c\}$ and the leader's sets
satisfy $\gamma(X)\le\beta$. MKIP (multidimensional KIP) allows $m_u\ge1$ and
$m\ge1$. In BOBILib's OR set, variant~$-100$ has one follower row ($m=1$) and
several leader rows; variants $-0$ and $-50$ have several follower rows. The
\emph{subset-sum class} of KIP is $p=w$; then $\Phi(X)$ is the largest subset
sum of $\{w_j: j\notin X\}$ that does not exceed $c$.

\paragraph{Clique and network interdiction.}
For edge clique interdiction, $G=(V,E)$ is a simple graph with $n=|V|$ and
$m=|E|$; the leader deletes $D\subseteq E$ with $|D|\le k$, and the follower
takes a maximum clique of $G-D=(V,E\setminus D)$, of size $\omega(G-D)$. In the
attack-protect network game the leader attacks at most $\Phi$ nodes and the
follower then protects at most $\Lambda$ nodes (Section~\ref{app:network}; there
$\Phi$ is the attack budget, as in BOBILib's row name, not a value function).

\paragraph{Letters with two meanings.} Following the paper, a few letters
are reused; the second meaning is local to one section, and the context
always says which is meant: $c$ is the leader's cost vector
in~\eqref{eq:bl} and the follower's capacity in packing interdiction; $D$ is
the follower's matrix in~\eqref{eq:bl} and the deleted edges in
Section~\ref{app:clique}; $G$ is the leader's budget matrix and the graph of
Section~\ref{app:clique}; $m$ is the number of follower rows in
Section~\ref{app:mkip} and the number of edges in Section~\ref{app:clique};
$k$ is an index or a count elsewhere (the depth of a search node in
Section~\ref{app:kip}, a row in Section~\ref{app:mkip}) and the deletion
budget in Section~\ref{app:clique}; $w$ denotes follower weights, the width of a
band table (Section~\ref{app:kip}) and node weights
(Section~\ref{app:network}); $\Phi$ is the follower's value in packing
interdiction and the attack budget in Section~\ref{app:network}. The value of
the online game is $\Omega$ (Sections~\ref{app:kip} and~\ref{app:mkip}), and
$\omega$ is the clique number only; a recession direction of the follower is
$r$ (Section~\ref{app:general}), not to be confused with the cost vector~$d$.

\subsection{Tolerances of the two checkers}\label{app:tolerances}

BOBILib's feasibility checker (described in Section~5.2 of the BOBILib
paper by Th\"urauf, Kleinert, Ljubi\'c, Ralphs and Schmidt,\cite{thurauf2026bobilib}
and distributed with the library as a script) accepts a
solution file in three steps: (i)~the HPR is feasible; (ii)~the file's point
passes the MIPLIB~2017 solution checker\cite{gleixner2021miplib} (version~1.0.3) on the HPR with
tolerance $10^{-4}$ for the linear rows, scaled by the size of the row (a relative test), and
$10^{-4}$ for integrality, and the
stated leader objective matches; (iii)~the follower's problem with the
leader's variables fixed to the file's values is solved to optimality (by
Gurobi with feasibility tolerance $10^{-6}$), giving $\varphi_{\mathrm{check}}$,
and the file's follower value $\varphi_{\mathrm{sol}}$ must satisfy
\begin{equation}\label{eq:bobilib-tol}
  \frac{|\varphi_{\mathrm{sol}}-\varphi_{\mathrm{check}}|}{10^{-10}+|\varphi_{\mathrm{check}}|}\le 10^{-4}.
\end{equation}
(The paper prints the denominator without the absolute value; the script uses
$|\varphi_{\mathrm{check}}|$. The script also contains an optimistic tie-break
test, which its driver does not call.)
Our verifier \texttt{bobi verify} checks (a)~every row, bound and integrality
requirement of the HPR at the file's point with absolute tolerance $10^{-6}$;
(b)~the stated leader objective; (c)~the follower's problem with the leader
fixed, re-solved by an exact knapsack oracle when there is one follower row
and by a HiGHS\cite{huangfu2018parallelizing} MILP otherwise, against test~\eqref{eq:bobilib-tol}, where
differences of at most $10^{-9}$ count as equal. A MILP stopped at a time
limit gives a verdict only when its incumbent is strictly better than the
file's value by more than~\eqref{eq:bobilib-tol} allows (``not optimal'') or
its proven bound passes~\eqref{eq:bobilib-tol} (``optimal''); otherwise the
verdict is ``unproven''. The follower problem contains only the follower's
rows, as in BOBILib's checker; coupling rows constrain the leader. Our
re-certification \texttt{bobi certify} (used for the points with coupling
rows) adds the optimistic tie-break test, and both of its optimality tests
use proven dual bounds and the relative tolerance $10^{-5}$.

The \emph{tight test} \texttt{bobi verify -{}-tight} keeps (a) and (b) but
replaces the relative test~\eqref{eq:bobilib-tol} in (c) by
\begin{equation}\label{eq:tight-tol}
  \varphi_{\mathrm{sol}}-\varphi\le\max\bigl(10^{-6},\,10^{-9}|\varphi|\bigr)
  \quad\text{for every follower optimum $\varphi$ the solver's proof allows,}
\end{equation}
with the follower solved to zero relative gap and absolute gap $10^{-9}$ (for an
integer-valued follower objective the bound is rounded up). A file that passes
BOBILib's test but not~\eqref{eq:tight-tol} uses BOBILib's tolerance on the
follower's value; it is a benchmark-accepted point (Section~\ref{app:classes}).
All \NumTightFiles{} solution files of the data package pass the tight test
(\texttt{results/notes/tight\_verify.md}). \texttt{bobi certify -{}-tight} applies
the same bound with proven dual bounds and also bounds the tie-break: the
leader's objective within $10^{-6}+10^{-9}|F|$ of the proven optimistic
tie-break value.

``Optimal'' in this appendix means: a bilevel feasible point in the sense of
Definition~\ref{def:bilevel-feasible} (checked within the tolerances above)
whose value equals a proven lower bound. How rigorous the point and the bound
are is given by their certificate class (Section~\ref{app:classes}): exact,
solver-certified or benchmark-accepted.
